\documentclass[11pt,a4paper]{article}
\usepackage[T1]{fontenc}
\usepackage{lmodern,amsmath,amssymb,amsthm,mathtools,booktabs,microtype}
\usepackage[margin=26mm]{geometry}
\usepackage[colorlinks=true,linkcolor=blue,citecolor=blue,urlcolor=blue]{hyperref}
\setlength{\emergencystretch}{2em}
\newtheorem{theorem}{Theorem}[section]
\newtheorem{lemma}[theorem]{Lemma}
\newtheorem{proposition}[theorem]{Proposition}
\newcommand{\Q}{\mathbb Q}\newcommand{\Z}{\mathbb Z}
\DeclareMathOperator{\Tr}{Tr}\DeclareMathOperator{\Nm}{Nm}
\title{A quartic in 68 variables from translated secants}
\author{}\date{17 November 2026}
\begin{document}\maketitle
\begin{abstract}
Under ordinary rational $abc$ and the stated elliptic, prime-pattern,
trace-height and simultaneous-twist premises, we construct a defining
quartic for $\Z$ in $\Q$ with 68 total variables. Four quadrics in six
field coordinates describe two points with their reciprocal abscissas
and inverse formal parameters. Translation by a fixed two-torsion point
makes the abscissa of their difference a quadratic output, whose rational
half-trace needs only one additional coordinate. The chart is complete
across tangencies and every finite two-torsion difference. The complete
polynomial is a sum of 56 primitive integral quadratic squares, has 1078
coordinate monomials, and has integer content one. Its exact coefficient
height is $B^2+1$, where the fixed odd shift $B\ge81$ must be sufficiently
large. Its sufficient threshold is ineffective. No new guard or
arithmetic premise is required; zero has an explicit source witness.
\end{abstract}

\section{Statement and arithmetic premises}
Every variable count includes the input. For a reduced rational $u/v$,
$v>0$, write $h(u/v)=\log\max(|u|,v)$ and $h(0)=0$. Fix
\[
 K=\Q(\rho),\qquad \rho^2=2,\qquad D=11-6\rho,\qquad
 E:y^2=x^3-Dx,\qquad R_0=(1+\rho,2\rho).
\]
The arithmetic premises are those stated in~\cite{Arithmetic}: the
group description $E(K)=\Z R_0+E[2](K)$ with $R_0$ non-torsion;
the simultaneous-prime theorem for fixed contact data, with arbitrary
finite avoidance, prescribed Gaussian congruences, and local and real
approximation; the all-point trace-index bound~\cite{Trace}; and the
simultaneous congruent-twist statement in Section~\ref{sec:pole}.
We also assume ordinary rational $abc$. Write $m(Q)\in\Z$ for the
integer index defined by $Q=[m(Q)]R_0+T_0$ with $T_0\in E[2](K)$.
The algebraic charts below, including their exceptional strata,
do not require these arithmetic premises.

\begin{theorem}\label{thm:main}
There is a computable family
\[
 Q_B\in\Z[X,Y_1,\ldots,Y_{67}],\qquad
 B\ge81\text{ a fixed odd integer},
\]
whose members have degree four, integer content one, 1078 nonzero
coordinate monomials, and coefficient height
\[
                         H(Q_B)=B^2+1.
\]
Each member is a sum of 56 primitive integral quadratic squares.
Under the stated premises there is a fixed threshold $B_*$ such that
every odd $B\ge B_*$ gives rational input projection exactly $\Z$.
The shift $B$ is a coefficient and is not a quantified coordinate.
No numerical sufficient value of $B_*$ is asserted.
\end{theorem}

The exact coefficient specification is given in
\nolinkurl{computation/quartic/abc68_v65}. The equations and fixed
invertible row operations below specify that family completely.
The test value $B=81$ is not asserted to be a defining member.
The exact coefficient certificates verify finite identities and
coefficients, not every external arithmetic premise.

\section{The reference and translated-secant charts}
Every field coordinate denotes two rational coordinates in the basis
$1,\rho$. A bar denotes the nontrivial automorphism of $K/\Q$.
For a point with finite nonzero abscissa put
\[
                    f(Q)=1/x(Q),\qquad S_Q=y(Q)/x(Q).
\]
Use three reference field coordinates $u_H,w_H,v_H$. Put
\[
 a_0=1+\rho,\qquad c_0=4-2\rho,\qquad
 \delta_0=-23+17\rho.
\]
Then $R_0=(a_0,a_0c_0)$, $a_0-c_0^2=\delta_0$ and
$a_0\delta_0=D$. Impose the two field quadrics
\begin{equation}\label{eq:reference}
 a_0u_Hw_H=1,\qquad
 v_H\bigl(v_H-a_0(u_H+w_H)-2c_0\bigr)=\delta_0.
\end{equation}
The complete three-coordinate reference chart
\cite[Theorem~1.1 and Section~4]{Reference} identifies their
solutions exactly with
\[
 H\in E(K)\setminus\{O,(0,0),R_0,R_0+(0,0)\}.
\]
Its affine outputs are
\begin{equation}\label{eq:referenceoutputs}
 S=v_H-c_0-a_0u_H=S_H,\qquad
 T_H=a_0w_H-v_H+c_0=S_{H-R_0}.
\end{equation}
All three reference coordinates are units on the entire chart:
\[
 u_H^{-1}=a_0w_H,\quad w_H^{-1}=a_0u_H,\quad
 v_H^{-1}=\frac{v_H-a_0(u_H+w_H)-2c_0}{\delta_0}.
\]
The reciprocal abscissas are $f(H)=u_H/v_H$ and
$f(H-R_0)=w_H/v_H$. These are explanatory rational functions,
not additional coordinates. No inverse witness or extra guard
is required, including on the tangent and permitted torsion fibers.

Use rational coordinates $X,\lambda,b_2,b_4,t_1$ and set $\beta=X$.
The two power equations are
\begin{equation}\label{eq:powers}
                         b_2=X^2,\qquad b_4=b_2^2.
\end{equation}
Use six field coordinates $s,t,F,G,h,d$. The symbols $h,d$ in
this chart are independent of the reference point and its integer index.
Put $\pi=h+s+t$, an affine expression rather than a coordinate, and impose
the four field quadrics
\begin{equation}\label{eq:pair}
\begin{aligned}
 d&=\pi F,&
 D(F+G)&=(s-t)h-t^2,\\
 d&=Gh,&
 b_4&=d(s-t)-DFG.
\end{aligned}
\end{equation}
Define the field-valued output
\begin{equation}\label{eq:differenceoutput}
                            a=h\pi.
\end{equation}
The field element $a$ is not retained as two coordinates.

\begin{proposition}\label{prop:pair}
Suppose $\beta\ne0$. Equations~\eqref{eq:pair} describe exactly
ordered pairs $U,V$ on $E$ for which $U,V$ are finite with nonzero
abscissas and $U\ne V$. Their outputs satisfy
\[
 \begin{gathered}
 F=\beta^2f(U),\qquad G=\beta^2f(V),\qquad
 s=\beta S_U,\qquad t=\beta S_V,\\
                      a=\beta^2x(U-V).
 \end{gathered}
\]
Every tangent fiber is included, and the finite difference $U-V$ may
be any two-torsion point, including $(0,0)$. Neither $h,d,a$ nor either
ordinate is required to be nonzero. Over the field $K$, for fixed
$\beta\ne0$ the four quadrics define a smooth surface of dimension
two, and their Jacobian has rank four throughout it.
\end{proposition}
\begin{proof}
The equations give the exact identities
\begin{equation}\label{eq:pointidentities}
                  DF^2+s^2F=b_4,\qquad DG^2+t^2G=b_4.
\end{equation}
Indeed multiply the second row of~\eqref{eq:pair} by $F$ or $G$
and use the other rows. Since $b_4=\beta^4\ne0$, both $F,G$ are
units. On $E':y^2=x^3-D'x$, with $D'=D\beta^4$, recover
\begin{equation}\label{eq:recovery}
              U'=(\beta^4/F,s\beta^4/F),\qquad
              V'=(\beta^4/G,t\beta^4/G).
\end{equation}
Their curve equations are~\eqref{eq:pointidentities}.
If $U'=V'$, then $F=G$ and $s=t$. The two expressions for $d$
would give $s+t=0$, hence $s=t=0$. The second equation of
\eqref{eq:pair} would then give $2DF=0$, a contradiction.
Thus $U'\ne V'$.

Let $T_0=(0,0)$ and put $\mu=h+t$. Translation by $T_0$ gives
\[
         \widetilde U=U'+T_0=(-DF,DsF),\qquad
         \widetilde V=V'+T_0=(-DG,DtG).
\]
The line $y=\mu x+Dd$ passes through $\widetilde U$ and
$-\widetilde V$. The exact polynomial identity
\begin{equation}\label{eq:linefactor}
 Z^3-D'Z-(\mu Z+Dd)^2=(Z+DF)(Z+DG)(Z-a)
\end{equation}
follows from~\eqref{eq:pair}. Its third intersection point has
abscissa $a$, so the chord-and-tangent law, with intersection
multiplicities, gives
\[
               U'-V'=\widetilde U-\widetilde V
                     =(a,-\mu a-Dd).
\]
If $a=0$, then $h\pi=0$ and the two expressions for $d$ give
$d=0$; the difference is exactly $T_0$. No division by $a$ or $d$
has been made. Scaling back by $(x,y)\mapsto(x/\beta^2,y/\beta^3)$
gives the asserted interpretations on $E$.

Conversely, scale any allowed pair $U,V$ to $E'$ and define
$F,G,s,t$ by~\eqref{eq:recovery}. They satisfy
\eqref{eq:pointidentities}. The pair $(F,s)$ differs from $(G,t)$.
On the two open sets covering this domain, define
\begin{equation}\label{eq:inverseh}
 h=\frac{(s+t)F}{G-F}\quad(G\ne F),\qquad
 h=\frac{D(F+G)+t^2}{s-t}\quad(s\ne t).
\end{equation}
The difference of the two curve identities gives
\[
       (G-F)\bigl(D(F+G)+t^2\bigr)=F(s^2-t^2),
\]
so the definitions agree on the overlap. Set $d=Gh$.
The four equations~\eqref{eq:pair} follow directly, and the
recovery formulas are inverse. In particular $h,d$ are unique.

At a tangent fiber $F=G$, the equations force $t=-s\ne0$.
The second expression in~\eqref{eq:inverseh} gives
\[
                         h=\frac{2DF+s^2}{2s}.
\]
This is the regular tangent value; its denominator is nonzero.
The line-divisor identity~\eqref{eq:linefactor} also covers triple
intersections and nonzero-abscissa two-torsion differences. If an
individual ordinate vanishes, then its reciprocal abscissa is nonzero
by~\eqref{eq:pointidentities}; no inverse ordinate is introduced.

For the dimension and rank statement, the plane curve
$DF^2+s^2F=\beta^4$ is smooth: its two derivatives could vanish
simultaneously only if $s=0$ and $F=0$, which is impossible.
Equations~\eqref{eq:inverseh} identify the chart with the complement
of the diagonal in the product of two such curves. More explicitly,
eliminate $d$ using $d=Gh$. On $G-F\ne0$, solve the first equation
for $h$; the remaining two equations are an invertible linear
combination of the two curve equations. On $s-t\ne0$, solve the
second equation for $h$ and use the same curve equations. These two
opens cover the chart, so the four defining quadrics have Jacobian
rank four everywhere. This proves the smooth isomorphism and covers
all exceptional fibers in its domain.
\end{proof}

Retain only one rational coordinate $a^{(0)}$ for the output
\eqref{eq:differenceoutput}. If $h=h^{(0)}+\rho h^{(1)}$ and
$\pi=\pi^{(0)}+\rho\pi^{(1)}$, impose the rational quadric
\begin{equation}\label{eq:halftrace}
               a^{(0)}=h^{(0)}\pi^{(0)}+2h^{(1)}\pi^{(1)}.
\end{equation}
Then $\Tr(a)=2a^{(0)}$. The other component of $a$ is uniquely
$h^{(0)}\pi^{(1)}+h^{(1)}\pi^{(0)}$ and has no separate coordinate
or equation. Thus the chart and its half-trace use thirteen rational
coordinates and nine quadratic rows. For fixed nonzero $\beta$,
their rational geometric dimension is four and their Jacobian rank
is nine.

Two further field coordinates $J,K$ satisfy
\begin{equation}\label{eq:helpers}
                  v_HJ=u_HF,\qquad v_HK=w_HG.
\end{equation}
The automatic reference units give
\begin{equation}\label{eq:helpervalues}
 J=\beta^2f(H)f(U),\qquad K=\beta^2f(H-R_0)f(V).
\end{equation}
All of $F,G,J,K$ are nonzero when $\beta\ne0$. Their equations
are quadratic and introduce no inverse witness or nonzero guard.

Use two field coordinates for the comparison squares,
\begin{equation}\label{eq:comparisons}
                         C_2=(s-S)^2,\qquad I_2=(t-T_H)^2.
\end{equation}
The rational scaled trace and argument gap are
\begin{equation}\label{eq:gap}
 T_{\rm sc}=\Tr(a),\qquad
 t_1-\lambda=B+b_4^2+(\Tr(a))^2.
\end{equation}
Here $T_{\rm sc}=2a^{(0)}$ is affine in the retained rational
half-trace coordinate. The last square is $4(a^{(0)})^2$. Put $t_0=\lambda$ and
\[
                    g=t_1-t_0=B+\beta^8+T_{\rm sc}^2.
\]

\section{The contact and pole equations}
For each argument $t_j$, $j=0,1$, use twelve rational coordinates
\[
 R_j,Z_j,w_j,\eta_j,b_j,e_j,u_j,p_j,v_j,q_j,p_{2,j},q_{2,j}
\]
and three field coordinates
$A_{1,j},A_{2,j},\Omega_j=\omega_{0,j}+\rho\omega_{1,j}$.
For a single block suppress $j$ and denote its argument by $t_*$.
Impose
\begin{equation}\label{eq:contactrows}
 R=t_*^2-2t_*,\qquad RZ+Z-R^2+35R=0,\qquad
                         (R+1)w=Z^2+24Z,
\end{equation}
and
\begin{align}
 wA_1&=w^2+wC_2+J^2,\label{eq:factor1}\\
 wA_2&=w^2+w\overline{I_2}+\overline K^2,\label{eq:factor2}\\
 w\Omega&=A_1A_2,\label{eq:product}\\
 \omega_0^2-2\omega_0\eta-2\omega_1^2&=0,
                    \qquad wb=\eta^2,\label{eq:tracepole}\\
 p_2&=p^2,\qquad e^2+bp_2=1,\label{eq:pole1}\\
 uv&=e(1-b),\label{eq:pole2}\\
 q_2&=q^2,\qquad u^2-v^2=bq_2.\label{eq:pole3}
\end{align}
Every displayed row is quadratic. The exact count is
\begin{center}
\begin{tabular}{lrr}
\toprule
Part & Rational coordinates & Rational quadratic rows\\
\midrule
Input, powers and arguments &5&3\\
Reference chart &6&4\\
Translated-secant chart and half-trace &13&9\\
Reciprocal product helpers &4&4\\
Comparison squares &4&4\\
Two contact and pole blocks &36&32\\
\midrule
Total &68&56\\
\bottomrule
\end{tabular}
\end{center}

\begin{lemma}\label{lem:nonzero}
Every rational solution with $X\ne0$ has $w_0w_1\ne0$.
\end{lemma}
\begin{proof}
The automatic reference units and~\eqref{eq:helpervalues} give $J\ne0$.
If one weight were zero, its first factor equation would give
$J^2=0$, a contradiction.
\end{proof}

\section{Local extraction with inverse parameter errors}
At a good odd place, every finite valuation of a reciprocal abscissa
$f(Q)$ is even. Set $v(0)=+\infty$.

\begin{lemma}\label{lem:factor}
Suppose $v(w)=\nu>0$ is odd, and every finite valuation of $J$ is
even. Then $\mathcal D=1+C^2/w+J^2/w^2$ is nonzero and has
valuation at most zero. If $v(\mathcal D)>-\nu/2$, then
\[
                         v(C)>\nu/4,\qquad v(J)>3\nu/4.
\]
\end{lemma}
\begin{proof}
After discarding any zero summand, the nonconstant summands have
valuations $2v(C)-\nu$ and $2v(J)-2\nu$. The first, when present,
is odd; the second is $2$ modulo four. Neither finite valuation
is zero, and if either is negative the least valuation occurs
only once. The assertions follow by comparing these valuations
with $-\nu/2$. A discarded summand has infinite valuation and
satisfies the corresponding inequality automatically.
\end{proof}

\begin{lemma}\label{lem:inverse}
Let $H,U$ have finite nonzero abscissas at a good odd place.
If
\[
 v(\beta)\le0,\qquad v(\beta S_U-S_H)>0,\qquad
                       v(\beta^2f(H)f(U))>0,
\]
then both points are formal at $O$. Both ordinate ratios are
nonzero, and $S_H/(\beta S_U)$ is a unit congruent to one.
\end{lemma}
\begin{proof}
The curve identity is $Df(Q)^2+S_Q^2f(Q)=1$.
An exact zero ordinate forces $v(f(Q))=0$. If $S_H=0$, positivity
of the product forces $v(f(U))>0$, whence $v(S_U)<0$ and the
error cannot have positive valuation. The case $S_U=0$ is symmetric.
If $v(S_H)>0$, the error inequality gives $v(\beta S_U)>0$;
both reciprocal valuations are then zero, contrary to the product
inequality. If $v(S_H)=0$, the same contradiction follows from
$v(S_U)=-v(\beta)\ge0$. Consequently, writing $e=v(\beta)$,
\[
                         h=v(S_H)=v(\beta S_U)<0.
\]
If $h-e=v(S_U)\ge0$, then $v(f(U))=0$ and $e\le h$.
The two possible reciprocal valuations at $H$ are $-2h$ and
$2h$, so the product valuation is at most zero. Hence $h-e<0$.
For a negative ordinate-ratio valuation $z$, the reciprocal
valuation is $-2z$ on the formal branch at $O$, and $2z$ on the
branch at $(0,0)$. The four possible product valuations are
\[
\begin{array}{c|cc}
 & U\text{ at }O & U\text{ at }(0,0)\\\hline
 H\text{ at }O &4(e-h)&0\\
 H\text{ at }(0,0)&4e&4h
\end{array}
\]
Only the first can be positive. Thus both points are formal at $O$.
The positive error divided by $\beta S_U$, whose valuation is
negative, proves the final congruence.
\end{proof}

Write $m=m(U-V)\in\Z$. The contact rows give
\[
 w(t_*)=\frac{\prod_{\theta\in\Theta}(t_*-\theta)}{(t_*-1)^6},
 \qquad \Theta=\{-5,-2,-1,0,2,3,4,7\}.
\]
Outside the fixed exceptional set, a consumed contact prime has
positive odd order $\nu=v_p(w)$.

\begin{proposition}\label{prop:local}
For $X\ne0$, every consumed contact prime $p$ makes $\beta$
$p$-integral and satisfies
\[
                            \beta(m-\beta)\equiv0\pmod p.
\]
\end{proposition}
\begin{proof}
The physical factors $\mathcal T_i=A_i/w$ are
\[
 \mathcal T_1=1+(s-S)^2/w+J^2/w^2,\qquad
 \mathcal T_2=1+\overline{(t-T_H)}^{\,2}/w+\overline K^2/w^2.
\]
Both reciprocal products have even valuations.
Lemma~\ref{lem:factor} makes each factor nonzero of valuation at
most zero at both places. Set $W=\mathcal T_1\mathcal T_2$.
Then $\Omega=wW\ne0$. If $\omega_0=0$, its trace equation would
give $\omega_1=0$, a contradiction. It follows that
\begin{equation}\label{eq:physical}
 \tau=\frac{\Nm W}{\Tr W}\ne0,\qquad
                         \eta=w\tau,\qquad b=w\tau^2\ne0.
\end{equation}
The native pole obstruction~\cite{Pole} excludes negative odd
valuation of $b$, so $v_p(\tau)>-\nu/2$. The reciprocal-trace
inequality
\[
 v_p(\tau)\le\max_{\mathfrak p\mid p}v_{\mathfrak p}(W)
\]
selects an embedding at which both factor valuations exceed
$-\nu/2$. The separation lemma therefore gives positive valuations
of $s-S,J,\overline{t-T_H},\overline K$ at that embedding.
This controls the pair $H,U$ at one embedding and the pair
$H-R_0,V$ at its conjugate embedding.

If $v_p(\beta)>0$, the asserted residual already vanishes.
Otherwise Lemma~\ref{lem:inverse} places both reference and
numerator pairs in their formal groups and shows that their
physical parameter ratios
\[
             z(U)/z(H)=S_H/S_U,\qquad
             z(V)/z(H-R_0)=S_{H-R_0}/S_V,\qquad z=-x/y,
\]
have valuation $v_p(\beta)$ and are congruent to $\beta$ after
division by a unit congruent to one. The reference indices are
$h_0=m(H)$ and $h_0-1$, at least one of which is a $p$-adic unit.
At that reference, the formal-logarithm ratio is integral because
the numerator index is an integer. This forces $v_p(\beta)\ge0$,
and hence $v_p(\beta)=0$.

Both physical ratios now reduce to $\beta$. The formal logarithm
is an isometry and agrees with its parameter to positive relative
order, so
\[
 m(U)\equiv h_0\beta,\qquad m(V)\equiv(h_0-1)\beta\pmod p.
\]
For a divisible reference index, the integral raw ratio forces
the numerator index to be divisible too, giving the same congruence.
To include global torsion labels, multiply the point relations by
two before applying the logarithm; two is a unit at the good odd
place. Subtracting gives $m\equiv\beta$. This argument also applies
at an inert good prime, where the two embeddings preserve the
same valuation. All exact zero-ordinate and mixed torsion branches
were included in Lemma~\ref{lem:inverse}.
\end{proof}

\section{Scaled trace and height soundness}
If $X=0$, the input is an integer. Assume $X\ne0$.
By Lemma~\ref{lem:nonzero}, neither argument is a contact.
Ordinary rational $abc$ gives the eight-contact estimate
\cite[Section~2]{Height}
\[
                    h(t_j)\le\frac27\log M(t_j)+C_0.
\]
The mixed-power estimate at exponent $101/100$, applied to
$\beta^8+T_{\rm sc}^2$, gives
\begin{equation}\label{eq:height}
 h(\beta)\le\frac{301}{388}h(g)+C_1,\qquad
 h(T_{\rm sc})\le\frac{301}{97}h(g)+C_2,\qquad h(g)\ge\log B.
\end{equation}
The contact estimate uses ordinary $abc$ and is not asserted from
the single weak exponent alone. Proposition~\ref{prop:pair} gives
$T_{\rm sc}=\beta^2\Tr x(A)$ on every allowed pair fiber. Hence
\[
              h(\Tr x(A))\le h(T_{\rm sc})+2h(\beta).
\]
For $A\ne(0,0)$, the all-point trace-index theorem
~\cite[Corollary~1.3]{Trace} supplies fixed constants for the index.
If $A=(0,0)$, then $m=0$ and $\Tr x(A)=0$, so the same bound
holds trivially after taking its additive constant nonnegative.
Thus on every chart fiber,

\begin{equation}\label{eq:indexheight}
              m^2\le C\bigl(h(T_{\rm sc})+2h(\beta)\bigr)+J_0.
\end{equation}
This includes zero traces and every finite torsion difference in
the pair chart.

Suppose $m\ne\beta$. Write $\beta=c_1/d_1$ in lowest terms,
with $c_1\ne0$, $d_1>0$, and put
\[
                         N=|c_1|\,|md_1-c_1|>0.
\]
Proposition~\ref{prop:local} gives $p\nmid d_1$ and $p\mid N$
at each consumed prime. Thus both $M(t_j)$ divide $N$, and
\[
 h(g)\le\frac47\log N+C_3,\qquad
               N\le(1+|m|)\exp(2h(\beta)).
\]
Combining these bounds with~\eqref{eq:height} yields
\[
 h(\beta)\le\frac{602}{679}h(\beta)
        +\frac{301}{679}\log(1+|m|)+C_4.
\]
Strict absorption bounds $h(\beta)$, then $h(g)$ and
$h(T_{\rm sc})$, by fixed multiples of $\log(1+|m|)+1$.
Equation~\eqref{eq:indexheight} then bounds $|m|$ and all these
heights, independently of $B$. Its additional $2h(\beta)$ term
changes only this final constant. The resulting bound contradicts
$h(g)\ge\log B$ for every sufficiently large fixed odd $B$.
Consequently $X=\beta=m\in\Z$, proving soundness.

\section{Prepared pole inputs greater than one}\label{sec:pole}
The retained simultaneous-twist statement says: if $\beta_0>1$
and $\theta\ge1$ are coprime positive squarefree integers, both
$1\pmod8$, and each is a square modulo every prime of the other,
then some positive squarefree $r$ makes
$E_r:y^2=x^3-r^2x$ and $E_{r\theta}$ have infinite-order rational
points with abscissas in $\beta_0\Q^{\times2}$.
This is the retained application of Jaswal's
Theorem~1.2~\cite{Jaswal}, with the precise scope stated in~\cite{Pole}.

\begin{lemma}\label{lem:largepole}
Assume that statement. Suppose $b>1$ has square reduced denominator,
$b\in\Q_2^{\times2}$, $v_2(b)\le-4$, and every prime dividing
the squarefree part of its numerator is $1\pmod4$.
Then
\[
 e^2+bp^2=1,\qquad uv=e(1-b),\qquad u^2-v^2=bq^2
\]
has rational witnesses with $upq\ne0$.
\end{lemma}
\begin{proof}
Write $b=A/B_0^2$ reduced, with $A,B_0>0$, and write
$b=\beta_0 z^2$, $b-1=\theta c^2$, with $\beta_0,\theta$
positive squarefree and $z,c\ne0$. They are coprime, because
$\gcd(A,A-B_0^2)=1$.
The assumptions imply $4\mid B_0$, $A$ odd and $A\equiv1\pmod8$.
Thus both squarefree parts are $1\pmod8$.
At a prime dividing $\beta_0$, the residue of $b-1$ is $-1$;
it is a square because that prime is $1\pmod4$.
At a prime dividing $\theta$, the residue of $b$ is one.
Neither kind of prime divides $B_0$. These observations give the
two required mutual square-residue conditions.

If $\beta_0=1$, take $e=0$, $p=1/z$, $u=z$, $v=0$, $q=1$.
Otherwise apply the simultaneous-twist statement. Scaling its second
point by $(x,y)\mapsto(c^2x,c^3y)$ gives an infinite-order point
on $E_{r(b-1)}=E_{r(1-b)}$. Write the two points as
$(ba_1^2,y_1)$ on $E_r$ and $(ba_2^2,y_2)$ on $E_{r(1-b)}$,
with $a_1a_2y_1y_2\ne0$. Put
\[
 e=\frac{r}{ba_1^2},\quad p=\frac{y_1}{b^2a_1^3},\quad
 u=\frac{a_2}{a_1},\quad v=\frac{e(1-b)}u,\quad
 q=\frac{y_2}{b^2a_1^3u^2}.
\]
The first curve equation gives $bp^2=1-e^2$.
The second gives $bq^2=u^2-e^2(1-b)^2/u^2=u^2-v^2$.
The middle equation is the definition of $v$.
All displayed divisions are by nonzero quantities. The sign of
$1-b$ is retained in $v$; its square in the elliptic equation
justifies the indicated curve equality.
\end{proof}

\section{Completeness and preparation}\label{sec:completeness}
Fix $X\in\Z\setminus\{0\}$ and an odd $B\ge81$. Before selecting
any prime, fix
\[
 \beta=X,\qquad A=[X]R_0,\qquad
 T_{\rm sc}=\beta^2\Tr x(A),\qquad g=B+\beta^8+T_{\rm sc}^2.
\]
The point $A$ is non-torsion, and both its abscissa and ordinate
are nonzero. The two contact lists $\Theta$ and $\Theta-g$ are
disjoint. Fix all denominator data, real neighborhoods, and finite
local avoidance before applying the simultaneous-prime theorem.
In particular, exclude numerator and denominator primes of the
nonzero abscissas and ordinates of $R_0$ and $A$, at both embeddings,
and all primes dividing $2X$. Any further fixed nonzero coordinates
needed for preparation may also be excluded before prime selection.

The arbitrary-avoidance application in~\cite[Section~2]{Primes}
supplies $\lambda$, arbitrarily deep prescribed denominators, and
compatible order congruences for an integer $h_0$. The selected
Gaussian generators have real part $9$ modulo $16$ and imaginary
part zero modulo $16$. The halving lemma~\cite[Section~5]{Half}
puts $R_0$ and its conjugate in the doubled subgroups of both
reductions. Thus the order moduli $\#E(\mathbf F_p)/4$ annihilate
$R_0$. The congruences $h_0\equiv0$ at the primary order and
$h_0\equiv1$ at the opposite order therefore suffice for
\[
 H=[h_0]R_0,\qquad U=[h_0X]R_0,\qquad V=[(h_0-1)X]R_0.
\]
This applies to odd and negative $X$ without a parity condition.

The compact-group progression argument permits fixed real
neighborhoods avoiding the finitely many poles of the point,
ordinate-ratio and chord functions on this progression. Apply it
to the finite family $Q,XQ,X(Q-R_0)$ with $Q=[h_0]R_0$.
The finitely many indices giving an identity or zero abscissa in
the reference or numerator points can also be omitted. Different
real components cause no sign restriction, since the detector
terms are squares. The local neighborhoods at the fixed
preparation primes are chosen at the same stage. In particular
$U-V=A$ is fixed, so the scaled trace and gap are unchanged.
Proposition~\ref{prop:pair} supplies the six pair coordinates,
including both nonzero reciprocal abscissas and the regular chord
value. Equation~\eqref{eq:halftrace} supplies the half-trace, and
\eqref{eq:helpers} supplies the two nonzero reciprocal products.

At a selected simple contact, $v_p(w)=1$ and $\beta$ is a unit.
At the primary embedding $Q=[h_0]R_0$ is formal, and formal
multiplication gives
\[
             z([X]Q)=Xz(Q)\bigl(1+O(z(Q)^4)\bigr).
\]
Since $S_Q=-1/z(Q)$, this implies
\[
                        \beta S_{[X]Q}-S_Q=O(z(Q)^3).
\]
The same expansion holds for $Q'=[h_0-1]R_0$ at its opposite
reference embedding. With the conjugation in the second factor,
both intended errors therefore have valuation at least three
at one embedding of the detector. Their squares divided by $w$
have valuation at least five. Each reciprocal abscissa there has
valuation at least two, so each squared product divided by $w^2$
has valuation at least six. Both intended factors are units
congruent to one.

At the other detector embedding the reference and numerator
points reduce to $\pm R_0$ and $\pm A$. The prescribed abscissa
avoidance makes both reciprocal abscissas units, so each product
term has valuation $-2$. The separation in
Lemma~\ref{lem:factor} prevents cancellation with the error term.
Thus each opposite factor has valuation at most $-2$.
The product $W$ is a unit at one embedding and has negative
valuation at the other. Its trace cannot undergo equal-order
cancellation, and $v_p(\tau)=0$. The physical pole input therefore
satisfies $v_p(b)=1$ at each selected prime.

The simultaneous-prime theorem in its arbitrary-avoidance form
\cite{Primes} selects every contact prime in the class
$1\pmod{8\cdot3\cdot17}$. Write $\lambda=u_0/L$, let
$g_0=0$, $g_1=g$, and absorb all fixed denominators into $L$.
Its positive prime values are
$p_{j,\theta}=u_0+(g_j-\theta)L$. The contact identity gives
\begin{equation}\label{eq:squareclass}
 w_j=
 \frac{\prod_{\theta\in\Theta}p_{j,\theta}}
      {L^2\bigl(u_0+(g_j-1)L\bigr)^6}.
\end{equation}
The denominator is a rational square. Hence the squareclass of
$b_j=w_j\tau_j^2$ is supported only on the selected primes, all
$1\pmod4$. Their valuations are positive; all other valuations
are even. The reduced denominator of $b_j$ is consequently a square.

At each fixed preparation prime, the point, chord and inverse-error
functions lie in the previously fixed bounded local neighborhoods.
Increasing the prescribed denominator depth makes $1/w_j$
arbitrarily small. The factors tend to one, $W_j$ tends to one,
and $\tau_j$ tends to $1/2$. At two, the deep denominator gives
$v_2(w_j)$ negative even and eventually $v_2(\tau_j)=-1$.
Thus $v_2(b_j)\le-4$. Formula~\eqref{eq:squareclass} also makes
$w_j$ a two-adic square, because all selected primes are
$1\pmod8$; therefore $b_j$ is a two-adic square.

On the fixed real compact neighborhoods, as both arguments
increase, $w_j\sim t_j^2$, the factors tend uniformly to one, and
\[
                            b_j\sim t_j^2/4.
\]
The prescribed sufficiently large real argument gives $b_j>1$.
Lemma~\ref{lem:largepole} applies to both actual physical inputs
with independent pole witnesses. Take
$A_i=w_j\mathcal T_i$, $\Omega=w_jW_j$, and $\eta=w_j\tau_j$;
all detector equations then hold. The pole witnesses and their
two squares satisfy the remaining rows. This proves completeness
for every nonzero integer input.

For $X=0$, choose a fixed allowed reference with $S\ne T_H$,
for example $H=2R_0$. For this point
$\Nm(S-T_H)=7199/328\ne0$. Put
\begin{equation}\label{eq:zero}
\begin{gathered}
 b_2=b_4=0,\qquad s=S,\qquad t=T_H,\\
 F=G=d=J=K=C_2=I_2=0,\qquad
 h=\frac{t^2}{s-t},\\
 a=h(h+s+t)=\frac{s^2t^2}{(s-t)^2},\qquad
 a^{(0)}=\frac12\Tr(a).
\end{gathered}
\end{equation}
The two product equations and the scale equation in~\eqref{eq:pair}
vanish identically. Its remaining row is $(s-t)h=t^2$, which
holds by definition. The half-trace and helper rows hold as well.
Put $c=-st/(s-t)$, so that $a=c^2$.
The fixed gap is now $B+(\Tr(c^2))^2$. Apply the same prime
theorem to these two fixed contact lists, with a sufficiently
large real argument and sufficiently deep two-adic denominators.
In each block set
\[
                A_1=A_2=w,\qquad\Omega=w,\qquad
                \eta=w/2,\qquad b=w/4.
\]
The detector rows hold, and~\eqref{eq:squareclass} gives the
same squareclass and square-denominator conclusions. The local
and real preparation gives a two-adic square with $v_2(b)\le-4$
and $b>1$. Lemma~\ref{lem:largepole} supplies the remaining
witnesses. Thus zero has a rational witness with no additional
coordinate or guard.

\section{Exact coefficient specification}
Take the rational components of the displayed field equations in
the basis $1,\rho$, clear denominators, divide by integer content,
and choose the sign making the first nonzero coefficient in the
stored monomial order positive. Write $\operatorname{prim}$ for
this normalization. The 68-entry coordinate order is listed in
\nolinkurl{computation/quartic/abc68_v65/quadratic-family.json}.

For the four primitive components of~\eqref{eq:reference}, in their
 displayed order, apply the fixed matrix
\[
 M=\begin{pmatrix}
 1&1&0&0\\
 -1&-2&0&0\\
 3&4&-1&1\\
 -7&-9&0&1
 \end{pmatrix}
\]
 and primitively normalize the resulting rows. This matrix has
 determinant $\pm1$. Together with the other 52 primitive component
 rows in the table order, this gives $r_1,\ldots,r_{56}$.
 Row $r_{37}$ is the first pole circle row and has constant
 coefficient one. Let $c_i$ be the coefficient of the constant
 coordinate monomial and parameter power $B^0$ in $r_i$. Define
\[
 q_{37}=r_{37},\qquad
 q_i=\operatorname{prim}(r_i-c_i r_{37})\quad(i\ne37),\qquad
                         Q_B=\sum_{i=1}^{56}q_i^2.
\]
 The term involving $B$ in the gap row is not included in $c_i$.
The stored basis file specifies all signs and primitive divisors.
These row operations are invertible over $\Q$. Since a sum of
rational squares vanishes exactly when each summand vanishes,
the rational zero set of $Q_B$ is precisely the full quadratic
source proved above.

Exact expansion gives 1078 nonzero coordinate monomials for every
$B\ge81$ and 1081 parameter-separated terms. The input occurs
only in $X^2-b_2$ and no constant elimination changes that row.
Thus the coefficient of $X^4$ is one, proving exact degree four
and integer content one for every integer $B$.
The constant coefficient is $B^2+1$, and the largest absolute
fixed coefficient is 3554. The other parameter-dependent
coefficients are $c_0+c_1B$, with $|c_1|\le8$ and no cancellation
for $B\ge81$. Their exact profiles satisfy
$|c_0+c_1B|\le B^2+1$ in that range. Hence
\[
                          H(Q_B)=\max\{3554,B^2+1\}.
\]
Since $B\ge81$, the fixed bound is already smaller than $B^2+1$.
Thus every member in the stated parameter range has height exactly
$B^2+1$. This algebraic bound does not give a sufficient defining
threshold.

The producer reconstructs the literal field equations and uses
the specified matrix, not a lattice-reduction search, to select
their basis. A separate checker reconstructs all 56 rows using
rational-pair sparse arithmetic, verifies the basis operations,
and independently expands their squares. Separate universal
polynomial identities verify the translated-secant chart, its point
and difference outputs, and the zero-input witness. Exhaustive
finite-field comparisons with an independently implemented elliptic
group law test tangent fibers, zero ordinates, and every two-torsion
difference, including $(0,0)$. Those finite comparisons supplement the all-field proof
in Proposition~\ref{prop:pair}; they do not replace it.

\begin{thebibliography}{99}
\bibitem{Arithmetic}\emph{A quartic in 74 variables with small integer coefficients},
\nolinkurl{papers/quartic_existence/quartic_abc74_20261116_v48.tex}.
\bibitem{Reference}\emph{A three-coordinate quadratic chart for an elliptic point and its translate},
\nolinkurl{papers/elliptic_charts/three_coordinate_reference_20261115_v40.tex}.
\bibitem{Trace}\emph{Trace heights on quadratic twists with $j=1728$},
\nolinkurl{papers/trace_heights/trace_on_all_points_20261113_v25.tex}.
\bibitem{Height}\emph{A quartic in 78 variables from a mixed-power height estimate},
\nolinkurl{papers/quartic_existence/quartic_abc78_20261115_v34.tex}.
\bibitem{Pole}\emph{Pole equations from simultaneous congruent-number twists},
\nolinkurl{papers/pole_parameters/congruent_twist_poles_20261114_v27.tex}.
\bibitem{Half}\emph{Heights of positive weighted elliptic traces},
\nolinkurl{papers/trace_heights/weighted_trace_20261114_v27.tex}.
\bibitem{Primes}\emph{Arithmetic dependencies of the construction with 29 variables},
\nolinkurl{papers/arithmetic_dependencies/arithmetic_dependencies_20261108_v8.tex},
Section~2.
\bibitem{Jaswal} M.~Jaswal, \emph{A proof that $\Z$ is Diophantine in $\Q$
through congruent-number twists}, Hexagon \texttt{2611.00331v2},
12 November 2026, Theorem~1.2, printed p.~4.
\url{https://hexagonmath.org/pdf/2611.00331v2}.
\end{thebibliography}
\end{document}
